\documentclass[../../../../main.tex]{subfiles}
\begin{document}

\begin{flushright}
\footnotesize\textit{【自動文字起こし・要確認】}
\end{flushright}

[解] $f(x) = x^2+7$

(1) $a^2+7$ が $2^n$ の倍数の時。$(a \in \mathbb{N}, n \in \mathbb{N} \ge 3) \cdots \text{①}$
\begin{align*}
f(a) &= a^2+7 \\
f(a+2^{n-1}) &= (a+2^{n-1})^2+7 = (a^2+7) + 2^n \cdot a + 2^{2n-2}
\end{align*}
$f(a)$ が $2^{n+1}$ の倍数の時、題意は成立。$f(a)$が $2^{n+1}$の倍数でない時、\\
①から $f(a) = a^2+7 = 2^n \cdot A \ (A \in \text{odd})$ とかける。
\begin{align*}
f(a+2^{n-1}) &= 2^n \cdot A + 2^n \cdot a + 2^{2n-2} \\
&= 2^n (A+a) + 2^{2n-2}
\end{align*}
ここで、$A, a \in \text{odd}$ から $A+a \equiv 0 \pmod 2$、又 $n \ge 3$ から $2n-2 \ge n+1$\\
だから各項はいずれも $2^{n+1}$ でわり切れる。つまり $f(a+2^{n-1})$ は $2^{n+1}$ でわり切れる\\
以上から示された。終

(2) $n$についてのキノウ法で示す。\\
$n=1,2,3$の時、$a_1=1, a_2=1, a_3=1$ とすれば成立。\\
$n=k \in \mathbb{N}$ での成立をかていすると、(1)から $f(a_k)$ 又は $f(a_k+2^{k-1})$ のいずれかは $2^{k+1}$ でわり切れるから、そのわり切れる方をとって $a_{k+1} = a_k \text{ or } a_k+2^{k-1}$ とすれば良く、$n=k+1$でも成立\\
以上から示された。

\end{document}
